\scanpage{341}
\mspage{80}
% archive sheet number: 340
\deleted{une autre \unsure{dans|de} $\underline{\mathrm{Ch}}_\Sigma(b,c)$, ce qui conduit (en regardant leurs images dans $\underline{\mathrm{Ty}}$) à se donner \unsure{encore}}
\deleted{\[
\varphi:\tau_1\longrightarrow\tau,\qquad\varphi':\tau'_1\mathrel{\unsure{-|\to}}\tau'\qquad\text{dans }\underline{\mathrm{Ty}},
\]}
\deleted{deux \inserted{couples d'}objets $(c,c'),(c_1,c'_1)$ de $\underline{\mathrm{Ch}}_{\Sigma,\tau}(a,b)\times\underline{\mathrm{Ch}}_{\Sigma,\tau'}(b,c)$ et de $\underline{\mathrm{Ch}}_{\Sigma,\tau_1}(a,b)\times\underline{\mathrm{Ch}}_{\Sigma,\tau'_1}(b,c)$ respectivement et des hom}
\deleted{\[
\left\{\begin{aligned}
c_1&\overset{\unsure{w|u}}{\mathrel{\unsure{-|\longrightarrow}}}\varphi^*(c) \qquad\text{\unsure{de}}\\
c'_1&\xrightarrow{w'}{\varphi'}^*(c')
\end{aligned}\right.
\]}
\deleted{et on voudrait en déduire une flèche de $\underline{\mathrm{Ch}}_{\Sigma,\deleted{\unsure{\tau\tau'}}}(a,c)$ au dessus de la flèche \deleted{$\varphi'*\varphi:\tau$}}
\deleted{\[
\varphi'*\varphi:\tau'_1*\tau_1\longrightarrow\tau'\cdot\tau.
\]}

\deleted{Mais voilà que}

Mais voilà que je me rappelle que ce fameux produit cartésien $p_1^*(\text{\unsure{fourbi}})\times p_2^*(\text{\unsure{machin}})$ n'est autre que le produit cartésien tout bête
\[
\underline{\mathrm{Ch}}_\Sigma(a,b)\times\underline{\mathrm{Ch}}_\Sigma(b,c),
\]
envoyé sur $\underline{\mathrm{Ty}}\times\underline{\mathrm{Ty}}$ par le produit cartésien des deux foncteurs de projection
\[
\begin{aligned}
p_{a,b}:{}&\underline{\mathrm{Ch}}_\Sigma(a,b)\longrightarrow\mathrm{Ty}\\
p_{b,c}:{}&\underline{\mathrm{Ch}}_\Sigma(b,c)\longrightarrow\mathrm{Ty},\qquad\text{d'où}
\end{aligned}
\]
\diagram{\begin{tikzcd}[ampersand replacement=\&]
\mathrm{Ch}_\Sigma(a,b)\times\underline{\mathrm{Ch}}_\Sigma(b,c)\arrow[d,"{p_{a,b}\times p_{b,c}}"]\\
\underline{\mathrm{Ty}}\times\underline{\mathrm{Ty}}
\end{tikzcd}}
